\begin{textsample}{NEG Dim 2 – al – Score 1.00 – al/al\_ef\_38.txt}  \label{ex:f2_neg_035}
9.
O COMPONENTE CURRICULAR ARTE, NO ENSINO FUNDAMENTAL ARTE No Ensino Fundamental, o componente curricular Arte está centrado nas linguagens artísticas : Artes Visuais, Dança, Música e Teatro.
Essas linguagens articulam saberes referentes a produtos e fenômenos artísticos e envolvem as práticas de criar, ler, produzir, construir, exteriorizar e refletir sobre formas artísticas.
A sensibilidade, a intuição, o pensamento, as emoções e as subjetividades se manifestam como formas de expressão, comunicação e interação no processo de aprendizagem em Arte.
Na BNCC ( 2017 ), a organização do componente Arte, para o Ensino Fundamental, está dividida em dois ciclos : 1º ao 5º ( Anos Iniciais ) e 6º ao 9º ( Anos Finais ), o que significa tempo maior que o anual para que os estudantes atinjam as expectativas de aprendizagem.
Convém respeitar os diferentes ritmos de aprendizagens, pois não existe uma hierarquia e nem uma ordem a ser seguida para o desenvolvimento do trabalho.
No âmbito didático, favorece o uso de diferentes abordagens e graus de complexidade, de acordo com o ciclo.
No Referencial Curricular de Alagoas, o componente Arte, para o Ensino Fundamental, continua sendo trabalhado em ciclos ( 1º ao 5º dos Anos Iniciais e 6º ao 9º dos Anos Finais ) conforme a BNCC, só que o desenvolvimento desse trabalho é feito ano a ano, de acordo com as fases de progressão proposta nesse componente, que contempla o desenvolvimento das competências e habilidades propostas para o componente curricular Arte com um olhar progressivo e \textbf{espiral}, ou seja, a progressão das aprendizagens não está proposta de forma linear, rígida ou cumulativa com relação a cada linguagem artística ou objeto de conhecimento.
Nesse sentido, é proposto um movimento, onde as novas experiências estejam relacionadas às anteriores e às posteriores na aprendizagem em Arte com a criação de Desdobramentos Didático Pedagógico, onde a cada ano, fases de progressão irão complementar e orientar com possibilidades o desenvolvimento das habilidades.
Os conteúdos devem ser constantemente retomados durante o trabalho pedagógico do professor, procurando oferecer uma apropriada construção do conhecimento dos estudantes acerca das práticas, possibilitando o aumento da flexibilidade na delimitação dos currículos e a adequação às realidades locais.
Para garantir a integração, equilíbrio e continuidade dos processos de aprendizagens das crianças, adolescentes, jovens e adultos, é importante o respeito às singularidades de cada sujeito e as relações que eles estabelecem com os conhecimentos, assim como a natureza das mediações em cada etapa e períodos de transições.
Com esse olhar, propõe -se o desenvolvimento, principalmente nos períodos de transições, de estratégias de acolhimento e adaptações para crianças e adolescentes para possibilitar que essa nova etapa seja construída com base no que o estudante sabe e é capaz de fazer, dando continuidade de forma progressiva às experiências vividas na Educação Infantil, Anos Iniciais e Anos Finais.

% matched lemmas: espiral
\end{textsample}
